\chapter{Experimental Evaluation}
\label{chap:evaluation}
% Backward-compatible label for older references.
\label{chap:experiments}

The previous chapters introduced the data protocol, the two learned
genomic representations, and the framework used to construct marker
rankings. This chapter brings those components together.

The organization follows the questions that motivated the thesis
rather than the order in which scripts were executed. The first
section asks whether the complete genomic representation supports
breed classification. The following sections progressively remove
information: first training individuals and then genomic markers. The
chapter then compares alternative ways of choosing those markers and
finally considers whether the main methodological observations can be
reproduced in an independent ovine domain.

All numbers reported in this chapter come from frozen protocol-v2
runs: development cross-validation over the three fixed folds, a
single locked-test evaluation for the final models, and ranking or
panel experiments whose selection rules were fixed before the results
were inspected.

\section{Breed Classification from the Complete Genomic Representation}
\label{sec:breed_classification_evaluation}

The first evaluation establishes the reference classification
problem. Both representation-learning approaches receive the same
ADAPTmap cohort, the same development folds, and the same fold-specific
train-fitted preprocessing. The difference begins only after the
shared genomic matrix has been constructed.

For the variational approach, development includes the comparison of
the predefined latent dimensions. The configuration is selected
according to mean Macro-F1 across all three folds, with Balanced
Accuracy and Accuracy used only as tie-breakers.

\begin{figure}[ht]
  \centering
  \includegraphics[width=0.78\textwidth]{fig_rq1_latentdim_selection}
  \caption{Development cross-validation metrics as a function of the
  variational latent dimension. The four candidate sizes are close
  (within $0.3$ Macro-F1 points); the pre-registered selection rule
  chooses $d=96$, which attains the highest mean Macro-F1.}
  \label{fig:rq1_latentdim_selection}
\end{figure}

For the contrastive approach, the representation-learning objective is
fixed and the downstream breed classifier is a $3$-nearest-neighbor
model fitted only on training embeddings.

The two models are discussed separately before being compared. The
interpretation accounts for their different objectives: the
variational model is explicitly supervised for breed prediction,
whereas the contrastive model learns its representation without breed
labels.

\begin{table}[ht]
\centering
\caption{Breed-classification performance of the two learned genomic
representations. Cross-validation results are reported as mean
$\pm$ standard deviation across the three development folds, while
the locked-test results correspond to the final models trained on the
complete development set. The Random Forest row is the classical
full-panel baseline computed on the same folds; it is reused as
$S_{\mathrm{full}}$ in the marker-efficiency experiment and was not
evaluated on the locked test.}
\label{tab:rq1_final_comparison}
\begin{tabular}{lccc}
\hline
Model & Macro-F1 & Balanced Accuracy & Accuracy \\
\hline
VMGP, development CV
& $0.9796 \pm 0.0032$
& $0.9794 \pm 0.0024$
& $0.9799 \pm 0.0032$ \\

Contrastive, development CV
& $0.9320 \pm 0.0070$
& $0.9427 \pm 0.0068$
& $0.9524 \pm 0.0055$ \\

Random Forest, development CV
& $0.7189 \pm 0.0054$
& $0.7151 \pm 0.0022$
& $0.8338 \pm 0.0022$ \\

VMGP, locked test
& $\mathbf{0.9703}$
& $\mathbf{0.9716}$
& $\mathbf{0.9765}$ \\

Contrastive, locked test
& $0.9496$
& $0.9574$
& $0.9671$ \\
\hline
\end{tabular}
\end{table}

\subsection{Final Evaluation on the Locked Test Set}
\label{sec:locked_test_evaluation}

Cross-validation answers a development question: which configuration
should be carried forward? The locked test has a different purpose.
It estimates how the already fixed models behave on individuals that
were not used for preprocessing design, model selection, early
stopping, or hyperparameter decisions.

For final training, a new preprocessor is fitted on all development
individuals. Each selected model is then trained for the fixed
$E_{\mathrm{final}}$ determined from the median best checkpoint epoch
across development folds. Only after both requested final models have
finished training is the locked test transformed with the frozen
development preprocessor.

For the variational model, breed logits are evaluated directly. For
the contrastive model, deterministic development embeddings are used
to fit the final $k=3$ nearest-neighbor classifier and the locked test
embeddings are then classified once.

The locked-test scores must be treated as final descriptive results.
They are not a second opportunity to select another latent dimension,
change the epoch budget, or revise preprocessing.

For the two final models, the preprocessor fitted on all development
individuals retained $46{,}441$ markers. The epoch budgets derived from
the median development best-checkpoint epoch were
$E_{\mathrm{final}}=168$ for the variational model and
$E_{\mathrm{final}}=1216$ for the contrastive model. After both final
training runs had completed, the locked test ($426$ individuals) was
transformed once with the frozen development preprocessor and evaluated
once; no retraining, epoch change, preprocessing revision or model
selection followed the test evaluation.

The locked-test results are consistent with development
cross-validation. The variational model moves from
$0.9796 \pm 0.0032$ to $0.9703$ (a decrease of $0.93$ points), while
the contrastive model improves from $0.9320 \pm 0.0070$ to $0.9496$
($+1.76$ points). Under the same development folds, the classical
full-panel Random Forest reaches $0.7189 \pm 0.0054$ Macro-F1, so the
learned representations clearly outperform a classical classifier on
the complete marker panel in cross-validation ($+0.26$ and $+0.21$
Macro-F1 respectively).

\begin{figure}[ht]
  \centering
  \includegraphics[width=\textwidth]{fig_rq1_confusion_locked}
  \caption{Locked-test confusion matrices (row-normalized) for the two
  learned representations. Tick labels are canonical class indices in
  the alphabetical breed order of Table~\ref{tab:adaptmap_breed_mapping};
  the weakest breed SEA is index $29$. For each model all but six of the
  $34$ breeds are recalled perfectly.}
  \label{fig:rq1_confusion_locked}
\end{figure}

\begin{figure}[ht]
  \centering
  \includegraphics[width=\textwidth]{fig_rq1_latent_spaces}
  \caption{Development embeddings of the two final goat models. Left:
  PCA of the VMGP posterior mean ($d=96$); the largest breeds form
  distinct regions while several closely related breeds overlap.
  Right: the contrastive $3$-dimensional embeddings projected with the
  Equal Earth map, showing the breed structure organized on the unit
  sphere despite the extreme bottleneck.}
  \label{fig:rq1_latent_spaces}
\end{figure}

\section{Learning with Fewer Individuals}
\label{sec:sample_scarcity_validation}

The complete-cohort experiment demonstrates that the learned genomic
representations achieve high breed-classification performance when the
entire development population is available. In practical applications,
however, reference populations are frequently constrained by genotyping
costs or by the limited numerical size of local and endangered breeds.
An important methodological question is therefore how classification
generalization behaves when the number of available training individuals
per breed is systematically restricted.

To examine this behavior under controlled conditions, the sample-efficiency
experiment keeps the VMGP-inspired model selected during development
($d=96$) strictly frozen in its architecture, loss weights, and optimization
protocol. As detailed in Section~\ref{sec:sample_size_protocol}, the analysis
is restricted to a cohort of 15 breeds that possess at least 30 training
individuals in each development fold. We evaluate training sizes of
$N \in \{5, 10, 15, 20, 25, 30\}$ individuals per breed, corresponding to
total training sets ranging from $75$ to $450$ animals.

Because all data-dependent preprocessing steps---marker missingness filtering,
minor allele frequency filtering, genotype-mode imputation, and LD
pruning---are refitted solely on the sampled training animals, the experiment
evaluates end-to-end sample scarcity rather than the robustness of the
neural network alone. Across the three development folds and the five nested
sampling replicates ($\mathcal{R}=\{42, 43, 44, 45, 46\}$), each value of
$N$ is evaluated through 15 empirical measurements on unchanged held-out
validation sets.

Table~\ref{tab:rq2_sample_efficiency_results} reports the mean and sample
standard deviation ($ddof=1$) of the evaluated metrics, along with the
proportion of performance retained relative to the reference condition
at $N=30$ ($S_{\mathrm{ref}} = 0.9692$).

\begin{table}[ht]
\centering
\caption{Breed-classification performance of the frozen VMGP model ($d=96$)
across different training sample sizes per breed ($N$) on the 15-breed cohort.
Values represent the mean $\pm$ standard deviation across 15 measurements
(3 development folds $\times$ 5 nested sampling replicates). Retained
performance is expressed relative to the reference condition at $N=30$
($S_{\mathrm{ref}} = 0.9692$).}
\label{tab:rq2_sample_efficiency_results}
\begin{tabular}{cccccc}
\hline
$N$ & Total Training Animals & Macro-F1 & Balanced Accuracy & Accuracy & Retained (\%) \\
\hline
5  & 75  & $0.9048 \pm 0.0337$ & $0.9094 \pm 0.0264$ & $0.9424 \pm 0.0139$ & $93.35\%$ \\
10 & 150 & $0.9549 \pm 0.0099$ & $0.9552 \pm 0.0101$ & $0.9686 \pm 0.0070$ & $98.53\%$ \\
15 & 225 & $0.9590 \pm 0.0100$ & $0.9595 \pm 0.0096$ & $0.9716 \pm 0.0064$ & $98.95\%$ \\
20 & 300 & $0.9647 \pm 0.0099$ & $0.9657 \pm 0.0107$ & $0.9742 \pm 0.0066$ & $99.53\%$ \\
25 & 375 & $0.9661 \pm 0.0109$ & $0.9671 \pm 0.0106$ & $0.9749 \pm 0.0064$ & $99.67\%$ \\
30 & 450 & $0.9692 \pm 0.0079$ & $0.9696 \pm 0.0080$ & $0.9769 \pm 0.0050$ & $100.00\%$ \\
\hline
\end{tabular}
\end{table}

\begin{figure}[ht]
  \centering
  \includegraphics[width=0.78\textwidth]{fig_rq2_learning_curve}
  \caption{Sample-efficiency learning curve for the frozen VMGP model
  on the 15-breed cohort: mean $\pm$ standard deviation across the
  $15$ measurements per sample size, the reference
  $S_{\mathrm{ref}}$, the pre-registered $98\%$ retention threshold and
  the resulting $N_{\min}=10$.}
  \label{fig:rq2_learning_curve}
\end{figure}

\subsection{Observations on the Learning Trajectory}

The empirical learning trajectory (Figure~\ref{fig:rq2_learning_curve})
displays two prominent characteristics across the evaluated range:

\begin{enumerate}
    \item \textbf{Transition from $N=5$ to $N=10$:}
    At $N=5$ (75 training animals in total across 15 classes), the model
    achieves a mean Macro-F1 of $0.9048$ with a standard deviation of
    $0.0337$. The dispersion across replicates is noticeably higher in this
    regime than at any other sample size, indicating that performance is
    sensitive to the specific composition of the sampled individuals.
    Increasing the training allocation to $N=10$ individuals per breed
    produces the largest single gain observed across the experiment, with
    mean Macro-F1 rising by $5.01$ percentage points to $0.9549$.
    Concurrently, the standard deviation decreases to $0.0099$, reflecting a
    substantially more consistent outcome across sampling replicates and folds.

    \item \textbf{Plateau behavior beyond $N=10$:}
    For sample sizes above $N=10$, additional training individuals yield
    progressively smaller performance gains. Increasing the training set
    from $N=10$ (150 animals) to the reference setting of $N=30$ (450 animals)
    increases mean Macro-F1 by $1.43$ percentage points ($0.9549$ to $0.9692$).
    The intermediate increments between $N=10$, $N=15$, $N=20$, $N=25$, and
    $N=30$ show modest adjustments in the range of $0.14$ to $0.57$ percentage
    points, while the variability remains stable with standard deviations
    close to $0.01$.
\end{enumerate}

\subsection{Identification of the Operational Minimum Sample Size}

As defined in Section~\ref{sec:minimum_requirements}, the operational
minimum sample size $N_{\min}$ is the smallest value of $N \in \mathcal{N}$
satisfying

\begin{equation}
    \overline{F_1^{\mathrm{macro}}}(N)
    \geq
    (1 - \epsilon_N) S_{\mathrm{ref}},
\end{equation}

where the reference score is $S_{\mathrm{ref}} = \overline{F_1^{\mathrm{macro}}}(30) = 0.9692$
and the tolerance parameter was frozen at $\epsilon_N = 0.02$ prior to the
analysis. This yields an operational threshold of

\begin{equation}
    0.98 \times 0.969223 = 0.949838.
\end{equation}

Evaluating the empirical results against this criterion:
\begin{itemize}
    \item At $N=5$, the mean Macro-F1 of $0.9048$ represents $93.35\%$ of
    the reference score, which falls below the $98\%$ threshold.
    \item At $N=10$, the mean Macro-F1 of $0.9549$ corresponds to $98.53\%$
    of the reference score, satisfying the retention requirement.
\end{itemize}

Under the pre-specified protocol, the operational minimum training sample
size is therefore identified as

\begin{equation}
    N_{\min} = 10.
\end{equation}

\subsection{Discussion and Methodological Limitations}

The observed plateau suggests that, within this experimental setup, most of
the predictive population structure captured by the pipeline can be acquired
with relatively few training animals per breed. However, interpreting these
findings requires several methodological cautions:

\begin{itemize}
    \item \textbf{Confounded sources of sample scarcity:}
    Because the entire pipeline is re-executed for each subsample, the observed
    trajectory reflects the combined impact of sample size on both genomic
    preprocessing (such as allele-frequency estimation and LD pruning) and
    the optimization of the variational autoencoder. The present design does
    not isolate which of these two stages is more sensitive to sample reduction
    at $N=5$.

    \item \textbf{Cohort and task specificity:}
    The threshold $N_{\min} = 10$ is an operational property of the evaluated
    setup. It is obtained on a cohort of 15 breeds selected specifically
    because they had sufficient animals available across all development
    folds. These breeds may have distinct genetic differentiation that eases
    classification compared to more closely related, recently derived, or
    heavily admixed populations.

    \item \textbf{Dependence on experimental parameters:}
    The value of $N_{\min}$ is coupled to the choice of $\epsilon_N = 0.02$,
    the 50k marker array density, and the specific architecture of the
    supervised variational model. It should not be interpreted as an absolute
    biological limit on the number of animals required to characterize a
    breed, but rather as an empirical baseline for the behavior of this
    pipeline under reduced data availability.
\end{itemize}

\section{Learning with Fewer Genomic Markers}
\label{sec:marker_efficiency_validation}

The next experiment changes the other side of the high-dimensional
problem. Instead of reducing the number of individuals, it reduces the
number of available SNPs.

Starting from a global ranking constructed inside the training
partition, nested panels are evaluated from very small subsets toward
the complete post-QC genomic representation. The central question is
whether most of the classification information is concentrated in a
relatively small portion of the high-density array.

A useful result is not necessarily a reduced panel that improves
classification. If a much smaller panel preserves nearly the same
performance, that already indicates substantial redundancy in the
original genomic representation.

Rankings are constructed inside each fold's training partition for the
four methods, converted into the frozen nested panels, and evaluated
with the downstream Random Forest on the corresponding validation
partition. The complete post-QC panel provides the reference
$S_{\mathrm{full}} = 0.7189 \pm 0.0054$ Macro-F1, and random panels of
the same sizes provide the chance baseline (five subsets per fold and
size). Figure~\ref{fig:rq3_panel_curves} shows the resulting curves and
Table~\ref{tab:rq3_selected_k} reports the aggregate values at selected
panel sizes.

Three regimes are visible. At very small panels ($K=50$) no method
exceeds $0.61$ Macro-F1 and all are far below the full panel. In the
intermediate range the two classical rankings separate from the neural
attributions: at $K=200$, $F_{ST}$ reaches $0.7356$ and Random Forest
$0.7277$, against $0.6523$ for SHAP and $0.6271$ for Integrated
Gradients. Beyond $K=500$ the curves plateau: Random Forest peaks at
$0.7664$ ($K=2000$) and $F_{ST}$ at $0.7602$ ($K=1000$), while SHAP and
Integrated Gradients keep improving slowly and remain below the
full-panel reference ($0.7309$ and $0.7072$ at $K=5000$).

Under the pre-registered retention rule ($\epsilon_P=0.02$, threshold
$0.98\,S_{\mathrm{full}}=0.7045$) the smallest qualifying panels are
$P_{\min}=200$ for $F_{ST}$ and Random Forest, $P_{\min}=500$ for SHAP
and $P_{\min}=5000$ for Integrated Gradients; with the less conservative
$\epsilon_P=0.05$ (threshold $0.6829$) Integrated Gradients qualifies at
$K=2000$. Random panels never satisfy the $98\%$ criterion (best mean
$0.6790$ at $K=5000$) and only approach the $95\%$ threshold at the
largest size, which confirms that the four rankings carry information
beyond chance.

Two observations deserve emphasis. First, the $F_{ST}$ panel with $200$
markers already matches the full-panel reference ($0.7356$ versus
$0.7189$): a panel containing less than $0.5\%$ of the retained markers
preserves the breed-discrimination signal under this evaluator, and the
Random Forest ranking behaves similarly. Second, while the ordering of
the methods is stable in the intermediate range
($F_{ST} \gtrsim$ Random Forest $>$ SHAP $>$ Integrated Gradients), the
margin of Integrated Gradients over the random baseline is small at
large panels and its $P_{\min}$ is sensitive to the tolerance: at
$K=2000$ its mean Macro-F1 ($0.7035$) lies only $0.0010$ below the
$\epsilon_P=0.02$ threshold.

\begin{figure}[ht]
  \centering
  \includegraphics[width=0.85\textwidth]{fig_rq3_panel_curves}
  \caption{Marker-efficiency curves per ranking family, random-panel
  band, full-panel reference $S_{\mathrm{full}}$ (dashed) and
  $P_{\min}$ markers for $\epsilon_P=0.02$ (stars).}
  \label{fig:rq3_panel_curves}
\end{figure}

\begin{table}[ht]
\centering
\caption{Marker-efficiency results at selected panel sizes: mean
Macro-F1 $\pm$ standard deviation across the three development folds
(IG: Integrated Gradients). The random row aggregates five random
subsets per fold ($15$ observations). The full-panel reference is
$S_{\mathrm{full}} = 0.7189 \pm 0.0054$.}
\label{tab:rq3_selected_k}
\small
\setlength{\tabcolsep}{4pt}
\begin{tabular}{lccccc}
\hline
$K$ & $F_{ST}$ & Random Forest & VMGP SHAP & Contrastive IG & Random \\
\hline
$50$ & $0.5960 \pm 0.0210$ & $0.6060 \pm 0.0311$ & $0.5084 \pm 0.0109$ & $0.5155 \pm 0.0213$ & $0.4410 \pm 0.0311$ \\
$200$ & $0.7356 \pm 0.0029$ & $0.7277 \pm 0.0149$ & $0.6523 \pm 0.0237$ & $0.6271 \pm 0.0087$ & $0.5769 \pm 0.0189$ \\
$500$ & $0.7495 \pm 0.0047$ & $0.7428 \pm 0.0044$ & $0.7145 \pm 0.0311$ & $0.6640 \pm 0.0128$ & $0.6125 \pm 0.0211$ \\
$2000$ & $0.7529 \pm 0.0076$ & $0.7664 \pm 0.0118$ & $0.7199 \pm 0.0115$ & $0.7035 \pm 0.0182$ & $0.6664 \pm 0.0166$ \\
$5000$ & $0.7581 \pm 0.0191$ & $0.7466 \pm 0.0172$ & $0.7309 \pm 0.0115$ & $0.7072 \pm 0.0068$ & $0.6790 \pm 0.0182$ \\
\hline
\end{tabular}
\end{table}

\section{Comparing Marker-Selection Strategies}
\label{sec:marker_ranking_comparison}

RQ4 asks whether learned-model explanations identify useful markers
relative to more classical references.

The comparison should place all ranking methods on equal footing. Each
method constructs a ranking from the same legitimate training
information. The ranking is then converted into the same sequence of
panel sizes, and the resulting panels are evaluated with the same
downstream classification protocol.

The purpose of this comparison is not to show that a deep attribution
method produces larger raw importance values than $F_{ST}$ or Random
Forest. The scores have different scales and meanings. What can be
compared meaningfully is the predictive utility and stability of the
reduced panels that each ranking produces.

RQ4 compares the four ranking families introduced in
Figure~\ref{fig:ranking_families} under the same protocol: identical
training partitions, identical nested panel sizes and the same
downstream evaluator. Their panel utility was quantified in
Section~\ref{sec:marker_efficiency_validation}; this section asks
whether the methods identify the same markers.

Averaged over the three folds, the Jaccard overlap between $F_{ST}$ and
Random Forest grows from $8.8\%$ ($K=50$) to $15.2\%$ ($K=200$) and
$17.8\%$ ($K=500$), while all remaining pairs stay below $2.2\%$
(Figure~\ref{fig:rq4_jaccard}). At $K=50$ the neural rankings share no
marker at all with the classical ones. Because the panels are tiny
compared with the $\sim 46{,}000$-marker space, these percentages must
be calibrated against chance: two independent random panels of the same
size would overlap by $0.05\%$, $0.22\%$ and $0.54\%$ at the three
sizes. The $F_{ST}$--Random Forest agreement is therefore roughly
$30$--$160$ times the chance level, while the agreement between learned
attributions and classical rankings remains within a few multiples of
chance. Considering fold $0$ for illustration, the $F_{ST}$--Random
Forest overlap at $K=500$ is extreme under a hypergeometric null
($p<10^{-180}$), whereas the single marker shared by $F_{ST}$ and SHAP
at $K=200$ is compatible with chance ($p\approx 0.58$).

The low overlap does not imply that the methods disagree about the
underlying biology. With strong linkage disequilibrium, several
correlated markers can act as interchangeable tags of the same
population signal, so different rankings may select distinct but
genomically linked markers. What the experiment does show is that the
identity of the selected markers is far from unique under this task,
and that the downstream utility of a ranking
(Table~\ref{tab:rq4_pmin}) is a more meaningful comparison than
score-level agreement.

Four mechanisms are consistent with the weaker performance of the
learned attributions as marker selectors. First, the contrastive
representation is an extreme three-dimensional bottleneck: many
distinct genotype profiles are mapped close to the same point on the
sphere, so per-marker relevance is fundamentally ambiguous. Second, the
contrastive model is trained to be invariant to aggressive
augmentations (allele flips and masking with probabilities up to
$0.99$), which explicitly discourages reliance on individual markers.
Third, the Integrated Gradients baseline --- the zero vector in the
four-channel one-hot space --- is not a valid genotype configuration,
and attributions computed from out-of-distribution baselines are known
to be unstable. Fourth, and most importantly, the attribution methods
explain their own neural models, whereas the downstream panel evaluator
is a Random Forest; the Random Forest importance ranking benefits from
the same inductive bias as the evaluator. The experiment therefore
compares transfer across model families rather than the intrinsic
correctness of the attributions, and a different downstream evaluator
could change the ranking of the methods.
Section~\ref{sec:rq4_robustness} tests this prediction explicitly and
confirms it.

A final methodological remark concerns the difference between
attribution and selection. The earlier notebook-based attribution
experiment ranked markers on the complete dataset and evaluated panels
by cross-validation on the same individuals; that protocol is
optimistic. The results reported here construct every ranking inside
the training partition and evaluate panels on untouched validation
individuals, which is the only way the comparison can support claims
about unseen animals.

\begin{figure}[ht]
  \centering
  \includegraphics[width=\textwidth]{fig_rq4_jaccard}
  \caption{Pairwise Jaccard overlap between the four ranking families
  at $K=50,200,500$ (mean over folds, percentages). Chance-level
  overlap is $0.05$--$0.54\%$.}
  \label{fig:rq4_jaccard}
\end{figure}

\begin{table}[ht]
\centering
\caption{Operational minimum panel size per ranking method under the
pre-registered retention tolerances. $P_{\min}$ is the smallest grid
value whose mean Macro-F1 reaches
$(1-\epsilon_P)\,S_{\mathrm{full}}$; ``---'' indicates that no grid
value qualified.}
\label{tab:rq4_pmin}
\begin{tabular}{lcccc}
\hline
Method & $P_{\min}$ ($\epsilon_P=0.02$) & $P_{\min}$ ($\epsilon_P=0.05$) & Best mean Macro-F1 & at $K$ \\
\hline
$F_{ST}$ & $200$ & $200$ & $0.7602$ & $1000$ \\
Random Forest & $200$ & $200$ & $0.7664$ & $2000$ \\
VMGP SHAP & $500$ & $500$ & $0.7309$ & $5000$ \\
Contrastive IG & $5000$ & $2000$ & $0.7072$ & $5000$ \\
Random & --- & --- & $0.6790$ & $5000$ \\
\hline
\end{tabular}
\end{table}

\section{Robustness of the Marker-Selection Comparison}
\label{sec:rq4_robustness}

The comparison of Section~\ref{sec:marker_ranking_comparison} keeps the
downstream evaluator fixed: every panel is scored by the same Random
Forest, whichever ranking family produced it. The mechanistic reading
proposed there --- in particular the possibility that the Random Forest
importance ranking benefits from the same inductive bias as the
evaluator --- predicts that the ordering of the families could change
under a different downstream model. This section tests that prediction
with four train-only analyses on the primary caprine cohort: an
evaluator-sensitivity study, a model-native evaluation in which the
explained architecture is retrained on the panel, an extended
classical baseline, and a valid-baseline study of the contrastive
attributions. All analyses reuse the frozen folds and the cached
protocol-v2 rankings; no locked-test record is read, and every model is
fitted on the fold training partition only.

\paragraph{Evaluator sensitivity.}
Three additional downstream classifiers are applied to the same panels
on the same development folds: logistic regression and a linear SVM on
standardised dosages, and gradient boosting, all with hyperparameters
and seeds fixed before the run. Because the frozen full-panel reference
$S_{\mathrm{full}} = 0.7189$ belongs to the Random Forest evaluator,
each alternative evaluator is given its own reference on the complete
post-QC panel: logistic regression reaches $0.9686$, the linear SVM
$0.9656$ and gradient boosting $0.8633$, against $0.7189$ for the
frozen Random Forest. The Random Forest is therefore by far the weakest
of the four on the full marker space --- an expected consequence of the
$p \gg n$ regime --- and its retention thresholds are correspondingly
the easiest to satisfy.

Figure~\ref{fig:rq4_evaluator_sensitivity} and
Table~\ref{tab:rq4_evaluator_pmin} summarise the result. Under the
frozen evaluator the ordering of
Section~\ref{sec:marker_ranking_comparison} is reproduced: $F_{ST}$ and
Random Forest importance are the most marker-efficient
($P_{\min}=200$ at both tolerances), followed by SHAP ($500$) and
Integrated Gradients ($2000$/$5000$). Under the three alternative
evaluators the VMGP SHAP ranking becomes the most marker-efficient
family: at $K=200$ it leads by $3.5$ points under logistic regression
($0.8971 \pm 0.0094$ versus $0.8625 \pm 0.0302$ for $F_{ST}$) and by
$5.0$ points under the linear SVM ($0.8324$ versus $0.7819$); at
$K=2000$ it is the best ranking under all three evaluators, and under
gradient boosting the $F_{ST}$, Random Forest and SHAP rankings are
indistinguishable at $K=200$ ($0.789$--$0.795$) while Integrated
Gradients remains about $8$ points behind. The associated $P_{\min}$
values move accordingly: SHAP reaches the $5\%$ tolerance with $500$
markers under every alternative evaluator, and it is the only family
that can do so with $500$ markers under the strict $2\%$ tolerance
(logistic regression); under gradient boosting only SHAP and Random
Forest importance qualify within the grid ($2000$), and under the
linear SVM no family reaches the strict tolerance at all (SHAP misses
it by $0.0004$ at $K=2000$). The ranking comparison is therefore
evaluator-dependent: the classical advantage measured in
Section~\ref{sec:marker_ranking_comparison} is specific to the frozen
Random Forest, which shares its inductive bias with the Random Forest
importance ranking.

\paragraph{Native retraining.}
The second analysis asks a different question: do the top-ranked
markers carry the signal for the model that the attribution actually
describes? Every panel is re-evaluated by retraining the same
architecture on the panel markers only, with the frozen RQ1
hyperparameters, and scoring the fold validation partition with the
breed head (VMGP) or with the $k=3$ nearest-neighbour classifier on the
learned embeddings (contrastive).
Figure~\ref{fig:rq4_native_retraining} and
Table~\ref{tab:rq4_native_pmin} report the outcome. The full-input
references are $0.9796 \pm 0.0039$ for the VMGP and $0.9320 \pm 0.0086$
for the contrastive model. Once the VMGP is retrained on the panel, the
SHAP ranking is again the most marker-efficient ($0.8864$ at $K=200$
and $0.9447$ at $K=500$, versus $0.8677$ and $0.9229$ for $F_{ST}$),
and it reaches the $5\%$ tolerance with $500$ markers while the
classical and contrastive rankings need $2000$; at the strict tolerance
$F_{ST}$ and SHAP are the only families that qualify (at $K=2000$),
with Random Forest importance and Integrated Gradients missing it by
less than $0.001$. Under the contrastive evaluator the panels remain
much more marker-hungry --- the model approaches its $0.9320$
full-input reference only in the thousands of markers --- but the SHAP
panel is still the best at $K=2000$ ($0.9119$ versus $0.8981$ for
$F_{ST}$, $0.8769$ for Random Forest and $0.8741$ for Integrated
Gradients), and only SHAP and $F_{ST}$ reach the $5\%$ tolerance within
the grid. These results support the transferability reading of RQ4: the
attributions identify markers that are genuinely informative for the
model they explain, and the frozen Random Forest penalises them because
it does not exploit that signal.

\paragraph{Extended classical association baseline.}
Finally, the classical family is widened with a univariate chi-square
test of independence between genotype and breed, computed on the
training partition only --- the association statistic that plays the
role of a GWAS-style test in a task whose outcome is the breed label.
Under the frozen Random Forest evaluator it is the most economical
ranking of all: with $100$ markers it reaches both tolerances
($\epsilon_P=0.02$ and $\epsilon_P=0.05$), whereas $F_{ST}$ and Random
Forest importance require $200$
(Table~\ref{tab:rq4_association}). Its top-marker lists overlap
strongly with $F_{ST}$ ($50.1\%$ at $K=50$, $52.9\%$ at $K=200$,
$60.4\%$ at $K=500$, against a chance level of $0.05$--$0.54\%$),
moderately with Random Forest importance ($12.5$--$20.0\%$) and
essentially not at all with the neural attributions (at most $2.4\%$).
The association test therefore closes the classical family ---
label-aware univariate statistics and tree importance agree with each
other and are the best choice under a Random Forest deployment --- and
it reinforces the identity-level conclusion of
Section~\ref{sec:marker_ranking_comparison}: the neural and the
classical rankings select largely disjoint marker sets.

\paragraph{Valid-baseline Integrated Gradients.}
The mechanistic discussion of
Section~\ref{sec:marker_ranking_comparison} suggested that the weakness
of the contrastive attribution could be an artefact of its reference:
the zero vector of the four-channel one-hot space does not correspond to
any genotype. The hypothesis is testable directly. Recomputing the
centroid-path IG with the settings of the frozen protocol reproduces
the cached ranking exactly in its top $500$ markers (Jaccard $100\%$ at
$K=50$, $200$ and $500$) but drifts beyond that depth, so the tail of
the ranking is itself unstable and the recomputed control reaches the
$5\%$ tolerance at $K=1000$ against $2000$ for the cached run. Two
extended variants are then compared under the identical recomputed
protocol: a centroid-path IG whose baseline is the mean one-hot profile
of the training partition --- a valid genotype reference inside the data
distribution --- and a probe-path IG that attributes the output of a
supervised MLP probe trained on the training embeddings, again from the
mean baseline. As shown in Figure~\ref{fig:rq4_valid_baseline_ig},
neither variant restores competitiveness: at the $5\%$ tolerance the
zero-baseline control qualifies at $K=1000$ while the valid baseline and
the probe path require $2000$; at the $2\%$ tolerance only the control
qualifies ($K=5000$), with the probe missing the threshold by $0.0001$
at that size. The two variants overlap each other by about $30\%$ and
remain almost disjoint from the frozen rankings (at most $1.8\%$
against a chance level of $0.54\%$ at $K=500$). The
out-of-distribution-baseline explanation is therefore not supported by
the data: under the frozen evaluator the contrastive attributions remain
the weakest learned selector even with a valid reference and a
supervised read-out, and their ranking is unstable in the tail.

\begin{figure}[ht]
  \centering
  \includegraphics[width=\textwidth]{fig_rq4_valid_baseline_ig}
  \caption{Valid-baseline Integrated Gradients under the frozen Random
  Forest evaluator. Left: mean Macro-F1 over the three development
  folds for the recomputed zero-baseline control and the two extended
  variants, with the frozen $F_{ST}$ and SHAP curves as references and
  the full-panel line (dotted); right: mean-over-folds Jaccard overlap
  with the frozen rankings at $K=500$ (chance $0.54\%$). Neither a valid
  baseline nor a supervised probe read-out improves the contrastive
  attribution.}
  \label{fig:rq4_valid_baseline_ig}
\end{figure}

\paragraph{Synthesis.}
Taken together, the analyses above qualify the answer to RQ4 rather
than overturning it. Under the frozen protocol --- a Random Forest
downstream evaluator and transferability as the criterion --- classical
univariate rankings are the most marker-efficient; the chi-square
association test is the strongest of them, and the neural attributions
are not competitive, with Integrated Gradients consistently the weakest
learned attribution; for Integrated Gradients neither a valid
training-mean baseline nor a supervised probe read-out restores
competitiveness, which rules out the out-of-distribution baseline as
the main explanation. Under a different deployment model the picture
inverts for SHAP: when the panels are scored by logistic regression, a
linear SVM or gradient boosting, or when the explained model itself is
retrained on the panel, SHAP-selected panels are the most
marker-efficient family. The practical implication is that a ranking
should be validated, and the marker panel chosen, against the
classifier that will actually be deployed; universal transfer of a
ranking across model families should not be assumed. These analyses are
extensions of the frozen protocol and carry their own limitations: they
cover the caprine cohort only, use a coarse panel grid, fix the
alternative classifiers to standard untuned hyperparameters, and
evaluate panels on the development folds without a locked-test
confirmation.

\begin{table}[ht]
\centering
\caption{Operational minimum panel size per ranking family under the
frozen Random Forest evaluator and under the alternative downstream
classifiers of the robustness analysis (caprine cohort). Cells report
$P_{\min}$ for $\epsilon_P=0.05$ / $\epsilon_P=0.02$; ``---'' means
that no grid value qualified. The alternative evaluators use
$K\in\{50,200,500,2000\}$ and their own full-panel references (bottom
row); the Random Forest column reproduces the frozen grid and
Table~\ref{tab:rq4_pmin}.}
\label{tab:rq4_evaluator_pmin}
\small
\setlength{\tabcolsep}{4pt}
\begin{tabular}{lcccc}
\hline
Method & RF (frozen) & Logistic reg. & Linear SVM & Grad. boosting \\
\hline
$F_{ST}$ & $200$/$200$ & $500$/$2000$ & $2000$/--- & $500$/--- \\
Random Forest & $200$/$200$ & $500$/$2000$ & $2000$/--- & $500$/$2000$ \\
VMGP SHAP & $500$/$500$ & $500$/$500$ & $500$/--- & $500$/$2000$ \\
Contrastive IG & $2000$/$5000$ & $500$/$2000$ & $2000$/--- & $2000$/--- \\
\hline
$S_{\mathrm{full}}$ & $0.7189$ & $0.9686$ & $0.9656$ & $0.8633$ \\
\hline
\end{tabular}
\end{table}

\begin{figure}[ht]
  \centering
  \includegraphics[width=\textwidth]{fig_rq4_evaluator_sensitivity}
  \caption{Evaluator sensitivity of the marker-selection comparison
  (caprine development folds, mean $\pm$ standard deviation over the
  three folds). Each panel is a downstream classifier with its own
  full-panel reference (dashed); the frozen Random Forest is the only
  evaluator under which the classical rankings dominate the learned
  attributions.}
  \label{fig:rq4_evaluator_sensitivity}
\end{figure}

\begin{table}[ht]
\centering
\caption{Operational minimum panel size per ranking family when the
explained architecture is retrained on the panel markers only
(caprine cohort; $K\in\{200,500,2000\}$). Cells report
$P_{\min}(\epsilon_P=0.05)/P_{\min}(\epsilon_P=0.02)$ against each
model's own full-input reference (bottom row).}
\label{tab:rq4_native_pmin}
\begin{tabular}{lcc}
\hline
Method & VMGP breed head & Contrastive KNN \\
\hline
$F_{ST}$ & $2000$/$2000$ & $2000$/--- \\
Random Forest & $2000$/--- & ---/--- \\
VMGP SHAP & $500$/$2000$ & $2000$/--- \\
Contrastive IG & $2000$/--- & ---/--- \\
\hline
$S_{\mathrm{full}}$ (full input) & $0.9796$ & $0.9320$ \\
\hline
\end{tabular}
\end{table}

\begin{figure}[ht]
  \centering
  \includegraphics[width=\textwidth]{fig_rq4_native_retraining}
  \caption{Model-native panel evaluation: the explained architecture
  (VMGP breed head, left; contrastive embeddings with $k=3$ nearest
  neighbours, right) is retrained from scratch on the panel markers
  only, with frozen RQ1 hyperparameters. The dashed line is the same
  model trained on the full post-QC panel. SHAP-selected panels are the
  most marker-efficient under both architectures.}
  \label{fig:rq4_native_retraining}
\end{figure}

\begin{table}[ht]
\centering
\caption{Extended chi-square association baseline under the frozen
Random Forest evaluator (caprine cohort). The association ranking
reaches $P_{\min}=100$ at both tolerances; the table reports its
mean-over-folds Jaccard overlap with the four frozen ranking families
at three panel sizes (chance level $0.05$--$0.54\%$).}
\label{tab:rq4_association}
\begin{tabular}{lcccc}
\hline
 & $F_{ST}$ & Random Forest & VMGP SHAP & Contrastive IG \\
\hline
$K=50$ & $50.1\%$ & $12.5\%$ & $0.0\%$ & $0.3\%$ \\
$K=200$ & $52.9\%$ & $17.6\%$ & $0.3\%$ & $0.9\%$ \\
$K=500$ & $60.4\%$ & $20.0\%$ & $0.8\%$ & $2.4\%$ \\
\hline
\end{tabular}
\end{table}

\section{Cross-Species Evaluation on Sheep HapMap}
\label{sec:cross_species_evaluation}

The final evaluation moves outside the primary caprine domain. An
independent ovine Sheep HapMap dataset is processed using the same
methodological principles but with a completely separate
species-specific pipeline.

This experiment should not compare goat and sheep raw scores as if the
two classification problems had identical difficulty. The number of
breeds, class balance, marker array, LD structure, and available
sample sizes may differ substantially.

The more meaningful comparison concerns the behavior of the method.
For example, one can ask whether compact representations remain
discriminative, whether performance responds similarly to sample
scarcity, or whether reduced panels retain useful information in both
species.

Within the ovine domain the pipeline is repeated independently: the
ovine split is frozen before any model is trained, every preprocessor
is fitted on ovine training individuals only, the caprine fold-specific
panels are never reused, and the same fractions and seeds produce a
development set of $1585$ individuals, a locked test of $278$
individuals and three development folds over the $28$ retained breeds
(Section~\ref{subsec:sheep_hapmap_dataset}). The same downstream
protocol is applied: three-fold development cross-validation and one
final training run per model, followed by a single locked-test
evaluation.

Table~\ref{tab:cross_species_rq1} and
Figure~\ref{fig:cross_species_rq1} summarize the classification
results.

\begin{table}[ht]
\centering
\caption{Breed classification in the two domains (Macro-F1). Goat
values come from the primary cohort, sheep values from the independent
ovine cohort. Cross-validation entries are mean $\pm$ standard
deviation over the three development folds; locked-test entries are
single evaluations of the final models. The goat Random Forest was not
evaluated on the locked test: its development-CV value is the
$S_{\mathrm{full}}$ reference of RQ3.}
\label{tab:cross_species_rq1}
\begin{tabular}{lcccc}
\hline
Method & Goat CV & Goat locked & Sheep CV & Sheep locked \\
\hline
Random Forest & $0.7189 \pm 0.0054$ & --- & $0.9168 \pm 0.0163$ & $0.9330$ \\
VMGP $d=96$ & $0.9796 \pm 0.0032$ & $0.9703$ & $0.9616 \pm 0.0133$ & $0.9555$ \\
Contrastive emb$=3$ & $0.9320 \pm 0.0070$ & $0.9496$ & $0.9160 \pm 0.0174$ & $0.9148$ \\
\hline
\end{tabular}
\end{table}

Three observations follow. First, both learned representations remain
strongly discriminative on an independent ovine cohort without any
species-specific hyperparameter tuning: the transferred VMGP reaches
$0.9555$ Macro-F1 on the ovine locked test and the contrastive model
$0.9148$. Second, the advantage of the learned representations over
the classical full-panel Random Forest is much smaller in the
validation domain. Under development cross-validation the VMGP exceeds
the Random Forest by $26.1$ points in goats but only $4.5$ points in
sheep, and on the ovine locked test the Random Forest ($0.9330$) lies
between the two learned models. Third, the self-supervised
contrastive representation does not outperform the classical baseline
on the ovine locked test ($0.9148$ versus $0.9330$), whereas on the
caprine cohort it did ($0.9496$ versus a development-CV Random Forest
reference of $0.7189$; the goat Random Forest was not evaluated on the
locked test, so that comparison remains indicative).

The cross-species evidence therefore supports the claim that compact
genomic representations generalize across species, but qualifies the
claim that they uniformly dominate classical baselines: the size of
the advantage is dataset-dependent, and on the ovine cohort the
classical full-panel classifier is a considerably stronger
competitor. As anticipated, the two cohorts differ in breed count,
sample sizes, marker panel and population structure, so the comparison
concerns the behaviour of the methods rather than a ranking of the two
datasets.

\begin{figure}[ht]
  \centering
  \includegraphics[width=\textwidth]{fig_cross_species_rq1}
  \caption{Cross-species comparison of RQ1. Left: development
  cross-validation Macro-F1 per method; right: locked-test Macro-F1
  where available. The gap between the learned representations and the
  classical Random Forest shrinks dramatically from the caprine to the
  ovine domain.}
  \label{fig:cross_species_rq1}
\end{figure}

\begin{figure}[ht]
  \centering
  \includegraphics[width=\textwidth]{fig_sheep_contrastive_embedding_space}
  \caption{Ovine contrastive embeddings on the unit sphere (validation
  folds) and their Equal Earth projection, colored by breed. Breed
  structure emerges without breed supervision despite the
  three-dimensional bottleneck.}
  \label{fig:sheep_embedding_space}
\end{figure}

% fig_sheep_confusion_locked.pdf and fig_sheep_vmgp_training_curves.pdf
% are available in assets/ for the appendix.

% ---- BEGIN AUTO RQ2-MINI ----
\paragraph{Ovine sample efficiency.}
The bounded mini-sweep (Section~\ref{sec:minimum_requirements}) evaluates
$N\in\{5, 10, 15\}$ individuals per breed with three subsampling seeds and
three development folds, fitting a fresh preprocessor for every run.
Retention is measured against the full-sample ovine VMGP baseline
($S_{\mathrm{full}} = 0.9616$). Across the 9 observations per sample
size, $N=5$: Macro-F1 $0.9273 \pm 0.0093$, retention $96.4 \pm 1.0\%$; $N=10$: Macro-F1 $0.9481 \pm 0.0075$, retention $98.6 \pm 0.8\%$; $N=15$: Macro-F1 $0.9531 \pm 0.0058$, retention $99.1 \pm 0.6\%$. At $N=10$ the ovine retention ($98.6\%$) already parallels the caprine result, where $N_{\min}=10$ retained $98.5\%$ of the reference. As in the caprine domain, most of the full-sample
performance is already acquired with very few individuals per breed and
the marginal gain beyond $N=10$ is small. Because the minimum
fold-training count per breed is $17$, the grid cannot be extended beyond
$N=15$ and the internal reference is $N=15$ itself; the ovine
$N_{\min}$ is therefore bounded by construction and is interpreted
through the shape of the retention curve rather than as an extended
minimum.

\begin{figure}[ht]
  \centering
  \includegraphics[width=0.80\textwidth]{fig_sheep_retention}
  \caption{Ovine bounded sample-efficiency mini-sweep: Macro-F1 retention
  (left) and absolute Macro-F1 (right) as a function of $N$, with error
  bars across the 9 fold-seed observations. The $98\%$ and $95\%$
  lines are the pre-registered retention thresholds and the dashed line is
  the full-sample baseline (0.9616).}
  \label{fig:sheep_retention}
\end{figure}
% ---- END AUTO RQ2-MINI ----

% ---- BEGIN AUTO RQ3-OVINE ----
\paragraph{Ovine marker efficiency and ranking comparison.}
The marker-reduction protocol of
Sections~\ref{sec:marker_efficiency_validation}--\ref{sec:marker_ranking_comparison}
is repeated inside the ovine cohort, with every ranking reconstructed
from ovine training individuals only and evaluated with the same
Random Forest downstream protocol. The full-array reference is the
ovine Random Forest, $S_{\mathrm{full}} = 0.9168 \pm 0.0163$, a much
stronger baseline than the caprine $0.7189 \pm 0.0054$.
Table~\ref{tab:cross_species_pmin} compares the operational minimum
panel sizes of the two domains; Figure~\ref{fig:cross_species_panels}
compares the retention curves, each normalised by its own
$S_{\mathrm{full}}$ so that the shapes of the curves, rather than the
absolute scores of two different classification problems, are
confronted.

Under $\epsilon_P=0.05$ the ovine Random Forest ranking is the most
economical ($P_{\min}=100$), $F_{ST}$ and VMGP SHAP tie at $200$,
contrastive Integrated Gradients requires $500$ and the random baseline
$1000$. At the stricter $\epsilon_P=0.02$ the ordering is Random Forest
$=200$, $F_{ST}=500$, SHAP $=1000$, and no IG panel qualifies: at
$K=5000$ its mean Macro-F1 ($0.8977$) remains below the $98\%$
threshold ($0.8985$). The family ordering of the caprine domain is
therefore reproduced in the validation domain --- classical rankings
and supervised attribution first, contrastive attribution last, random
panels worst --- with SHAP catching up with $F_{ST}$ at the $5\%$
tolerance.

Two quantitative differences qualify the replication. First, the ovine
panels saturate at larger sizes relative to their reference: in
caprines the $F_{ST}$ panel at $K=200$ already matches the full-panel
Random Forest ($0.7356$ versus $0.7189$), whereas in sheep the curves
converge to $S_{\mathrm{full}}$ only in the $500$--$1000$ range and
their best values ($F_{ST}$ $0.9252$ at $K=1000$, Random Forest
$0.9244$ at $K=2000$) exceed it by less than one point. This mirrors
the RQ1 cross-species result: the classical full-panel classifier is a
much stronger competitor in the ovine domain, so the claim that a small
panel matches the full array is weaker there. Second, the relative
penalty of the learned attributions narrows: at $\epsilon_P=0.05$ the
Integrated Gradients ranking needs ten times the Random Forest panel in
goats ($2000$ versus $200$) but five times in sheep ($500$ versus
$100$), and SHAP requires $2.5$ times the $F_{ST}$ panel in goats
($500$ versus $200$) but no additional markers in sheep.

At the identity level the Jaccard analysis replicates as well. Averaged
over folds, the $F_{ST}$--Random Forest agreement grows from $19.6\%$
($K=50$) to $21.2\%$ at both $K=200$ and $K=500$, against a chance
level of $0.06$--$0.62\%$; this classical agreement is even higher than
the caprine $8.8$--$17.8\%$. All remaining pairs stay below $2.8\%$, so
the learned attributions identify largely disjoint marker sets in both
species. As in the caprine analysis, low overlap does not imply
disagreement about the biology: under strong linkage disequilibrium,
distinct rankings can select interchangeable tags of the same
population signal. The consistent cross-species pattern is that
downstream utility orders the ranking families in the same way, while
score-level identity does not.

\begin{table}[ht]
\centering
\caption{Operational minimum panel size $P_{\min}$ per ranking family
in the two domains under the pre-registered tolerances. ``---''
indicates that no grid value qualified; the caprine columns reproduce
Table~\ref{tab:rq4_pmin}.}
\label{tab:cross_species_pmin}
\begin{tabular}{lcccc}
\hline
 & \multicolumn{2}{c}{Caprine (primary)} & \multicolumn{2}{c}{Ovine (validation)} \\
\cline{2-5}
Method & $\epsilon_P=0.02$ & $\epsilon_P=0.05$ & $\epsilon_P=0.02$ & $\epsilon_P=0.05$ \\
\hline
$F_{ST}$ & $200$ & $200$ & $500$ & $200$ \\
Random Forest & $200$ & $200$ & $200$ & $100$ \\
VMGP SHAP & $500$ & $500$ & $1000$ & $200$ \\
Contrastive IG & $5000$ & $2000$ & --- & $500$ \\
Random & --- & --- & $5000$ & $1000$ \\
\hline
\end{tabular}
\end{table}

\begin{figure}[ht]
  \centering
  \includegraphics[width=\textwidth]{fig_cross_species_panels}
  \caption{Marker efficiency in the two domains: mean Macro-F1 across
  the three development folds, normalised by each domain's own
  full-array reference $S_{\mathrm{full}}$; stars mark $P_{\min}$ at
  $\epsilon_P=0.05$ and the dotted line is the $95\%$ retention
  threshold. The ordering of the ranking families is preserved, while
  the ovine curves saturate closer to their (stronger) reference.}
  \label{fig:cross_species_panels}
\end{figure}
% ---- END AUTO RQ3-OVINE ----

\section{Overall Discussion}
\label{sec:overall_discussion}

The experiments in this chapter address different dimensions of the
same computational problem. Breed classification asks whether
population structure can be recovered from high-dimensional genotype
profiles. Sample-efficiency analysis asks how much training evidence
is required to learn that structure. Marker-efficiency analysis asks
how much of the original genomic input is necessary to preserve it.
Finally, marker-ranking comparison asks whether different notions of
feature relevance lead to equally useful compressed panels.

Across both domains the learned representations clearly outperform
the classical full-panel Random Forest in classification
(Tables~\ref{tab:rq1_final_comparison} and
\ref{tab:cross_species_rq1}): the advantage is $26.1$ Macro-F1 points
in caprines and $4.5$ points in ovines under development
cross-validation, and the supervised VMGP is consistently ahead of the
self-supervised contrastive model. This ordering is expected, since
the variational model is optimized for the same objective that is
later evaluated, whereas the contrastive model must recover breed
structure without labels. The result that matters is that the
self-supervised space remains strongly discriminative ($0.9496$ and
$0.9148$ Macro-F1 on the two locked tests), supporting the hypothesis
that population structure is a dominant, recoverable signal in these
genotype matrices.

The sample-efficiency experiments show a very early plateau. In the
caprine cohort the operational minimum was $N_{\min}=10$ with $98.5\%$
of the reference performance, and in the ovine cohort ten individuals
per breed already retained $98.6\%$ of the full-sample baseline. Two
caveats bound this conclusion: the ovine grid is truncated at $N=15$
by the minimum fold-training count of $17$ per breed, and all values
are conditional on the same pipeline, cohort composition and
architecture. The practical implication is nevertheless relevant for
numerically small breeds: a few tens of animals per breed can already
provide a useful reference panel.

Marker reduction is remarkably effective. With $200$ $F_{ST}$-selected
markers the downstream classifier matches the full post-QC panel
($0.7356$ versus $0.7189$ Macro-F1), and all four ranking families
outperform random panels beyond $K=200$. Panel sizes of this order
agree with the ancestry-informative marker literature, where a few
hundred markers are routinely sufficient for population assignment,
and recent cattle applications report accurate assignment with
panels of $192$ markers \cite{ahmad2024cattleassigner}. The caprine
curve also shows that the full panel is not an upper bound: a reduced
panel can match or slightly exceed it, consistent with noise dilution
in the $p\gg n$ regime. In the ovine cohort the same ordering is
reproduced under the $5\%$ tolerance with even smaller panels
($100$--$200$ markers versus $500$ for contrastive IG), although the
classical full-array reference is a much stronger competitor there, so
the ovine panels saturate near it rather than exceeding it. This
advantage is nevertheless conditional on the frozen evaluator and on
the protocol: panels were selected and evaluated inside the same
development folds, and no external cohort was used to validate them.

The ranking comparison (RQ4) provides the second main finding: under
the frozen Random Forest evaluator, neural attributions are not
competitive marker selectors for a different downstream model. $F_{ST}$
and Random Forest importance share a modest
but strongly above-chance fraction of their top markers ($8.8$--$17.8\%$
in caprines, $19.6$--$21.2\%$ in ovines), the extended chi-square
association baseline overlaps $F_{ST}$ by $50$--$60\%$,
while the neural rankings overlap each other and the classical ones by
at most a few percent, close to the chance baseline at small panel
sizes. This pattern is consistent with two documented phenomena:
panels selected by different criteria can have minimal marker overlap
even when they perform similarly \cite{soundararajan2016minimal},
because strong linkage disequilibrium makes several correlated markers
interchangeable tags of the same signal; and post-hoc attributions of
a model do not necessarily transfer to another model or to selection
decisions. The mechanistic analysis in
Section~\ref{sec:marker_ranking_comparison} identifies plausible
contributors: the extreme three-dimensional contrastive bottleneck,
the augmentation-invariance objective, an out-of-distribution
Integrated Gradients baseline, and the mismatch between the explained
neural models and the Random Forest evaluator. The robustness analyses
of Section~\ref{sec:rq4_robustness} show that this last mismatch is
consequential: when the panels are scored by logistic regression, a
linear SVM or gradient boosting, or when the explained model itself is
retrained on the panel, the SHAP-selected panels become the most
marker-efficient family, while the Random Forest evaluator favours the
rankings that share its inductive bias. The practical conclusion is
therefore twofold: attributions should be validated through downstream
utility rather than interpreted as feature importance per se, and the
evaluator used for that validation should match the model that will
actually be deployed.

Cross-species validation supports the generalization of the
representation pipeline but qualifies its superiority. Both models
transfer to the ovine cohort without species-specific tuning;
however, the classical full-panel Random Forest becomes a much
stronger competitor on sheep, and the self-supervised representation
no longer beats it on the locked test. This is consistent with
reports that classical and relationship-based methods remain
competitive for breed assignment depending on the dataset
\cite{wilmot2023genomic}, and reinforces the methodological reading:
the comparison concerns the behaviour of the methods, not a ranking of
the two species. The main limits of the study are the finite cohorts
($34$ and $28$ breeds, with tens of individuals per breed), a single
marker array per species, a single downstream evaluator, and reduced
panels that were not validated on external cohorts. The protocol,
ranking definitions and evaluation rules were fixed before the locked
tests were opened.
